%% 第 10 章：Fitting 分解
%% 本文件由 tools/md2tex.py 自动生成，请勿直接编辑。

\chapter{Fitting 分解}


\begin{jthm}{Fitting 分解定理}{r:11:1}
设：

\[
A:V\to V
\]

是有限维空间上的算子，且：

\[
n\ge\dim V.
\]

则：

\[
V
=
\ker A^n
\oplus
\operatorname{range}A^n.
\]
\end{jthm}

\section{Fitting 分解的意义}

空间被分为两部分：

\begin{jitem}
\item %
$\ker A^n$：反复作用 $A$ 后最终消失的部分；
\item %
$\operatorname{range}A^n$：经过很多次作用后仍然保留下来的部分。
\end{jitem}

当：

\[
A=T-\lambda I,
\]

第一部分就是 $\lambda$ 的广义特征空间。

\begin{jproof}{证明思路}
当 $n$ 足够大时，零空间链与值域链都已经稳定。

由秩—零化度定理：

\[
\dim\ker A^n+\dim\operatorname{range}A^n
=
\dim V.
\]

因此，只需证明二者交集为零。

若某个向量既在 $\ker A^n$ 中，又在 $\operatorname{range}A^n$ 中，则它可以写成 $A^n(u)$，同时又会被 $A^n$ 消掉。

这意味着 $u$ 落入更高次的零空间。由于零空间链已经稳定，$u$ 其实已经在 $\ker A^n$ 中，因此：

\[
A^n(u)=0.
\]

所以交集只能是零向量。
\end{jproof}
